\chapter{The Second Law and the Cost of Information}
\label{chap:second-law-cost-of-information}

The thread of this chapter is a single claim, and it is the chapter's spine: information is physical. The
second law is not, at bottom, a statement about disorder; it is a statement about the cost of bookkeeping.
And there is a sharper point inside it, one this part has been building toward: the instrument's own
measuring is a physical act with a thermodynamic price. To measure is not to look on from outside; it is to
join the causal order and pay into it. Every distinction the instrument draws, every bit it erases, every
boundary it must reconcile, is settled in the same currency as heat. Five effects build the case, from the
demon that seems to cheat the second law for free to the equilibrium that turns out to be actively held.

\section{Maxwell's Demon: sorting is entropy-producing}
\label{se:maxwell-demon}

The demon is the instrument in disguise. Set a partition across a box of gas with a gate the demon controls,
and let it watch the molecules, opening the gate only for the fast ones heading one way and the slow ones the
other. After a while the fast molecules are gathered on one side, the slow on the other --- the gas has been
sorted, hot from cold, its entropy lowered, and apparently for nothing, since opening a frictionless gate
costs no work. This is the paradox: the demon seems to break the second law by mere observation.

The resolution is that observation is never mere. The demon's measurement $M$ and the gas's evolution $U$ do
not commute; if $[M, U] = 0$ the demon's looking would leave the causal order untouched and could be slid past
the dynamics for free, but it does not, and the non-commutativity forces an entropy balance,
\begin{equation}
  \Delta S_{\text{gas}} + \Delta S_{\text{demon}} = k_{\mathrm{B}}\ln\!\left|\Omega_{\text{joint}}\right| > 0 .
\end{equation}
Entropy here is Boltzmann's count: the logarithm of the number of microscopic configurations a macroscopic
state allows,
\begin{equation}
  S = k_{\mathrm{B}} \ln \Omega,
\end{equation}
the formula on his tombstone, turning a count of arrangements $\Omega$ into a thermodynamic entropy. The
demon's sort lowers the gas's $\Omega$ and raises the demon's own, and the joint count
$|\Omega_{\text{joint}}|$ can only grow. The demon cannot read a molecule's speed without becoming part of the
gas's causal order, and each thing it
reads is a distinction it must record. That record is physical --- it lives in the demon's own memory --- and
the entropy the demon takes out of the gas reappears, and more, inside it: each refinement of its internal
partition, $P_n \to P_{n+1}$, adds a distinction to the world's total count.

Picture the demon's notebook as it works. Every molecule it sorts costs an entry --- \emph{fast, this side;
slow, that side} --- and the notebook fills, line by line, with exactly the distinctions the demon drew out of
the gas. The gas grows orderly; the notebook grows full. Nothing has been gotten for free: the order taken
from the gas has been paid for, distinction for distinction, in the demon's own record. The demon is the
instrument, and its notebook is the instrument's count --- and the count, as always, is physical.

The honest floor is a finite count obligation: the global count of distinct configurations rises, and the
demon's record is exactly the count it adds. The reading is the chapter's spine. Measurement is physical. The
demon does not cheat the second law, because observing is not free --- to measure is to join the system's
causal order and add to the world's count of distinctions. Every reading the instrument has taken in this
book, from the first mark to here, has been such an act; the demon simply makes its price visible.

\section{The Thermodynamic Cost of Erasure: Landauer's bound}
\label{se:landauer}

There is a loophole to close, and closing it is Landauer's. The demon's notebook is finite. Sooner or later
it fills, and to keep sorting the demon must erase it --- wipe the old entries to make room for new.
Measuring, it turns out, can be done for free; erasing cannot. Resetting a single bit --- dumping its recorded
distinction into the surroundings --- releases at least
\begin{equation}
  E_{\min} = k_{\mathrm{B}} T \ln 2
\end{equation}
of heat and raises entropy by at least $k_{\mathrm{B}}\ln 2$. This is Landauer's bound, and it is exactly the
size of the demon's debt. When the demon erases its full notebook to carry on, it pays out, bit by bit,
precisely the entropy its sorting had removed from the gas. The loophole closes: the sort lowered the gas's
entropy, the erasure that lets the sort continue raises it right back, and the second law is whole. The bill
came due not at the measurement but at the erasure.

Szilard sharpened the demon to a single molecule and made the exchange exact. Put one molecule in a box and
let the demon learn which half it is in --- a single bit of information. With that bit the demon can slide a
piston against the empty side and let the molecule's own pressure push it out, extracting work; for the
isothermal expansion of one molecule from half the box to the whole, the work is exactly $k_{\mathrm{B}} T
\ln 2$. One bit of information, converted into one bit's worth of work. But the demon now holds a used bit ---
which half the molecule was in --- and to run the engine again it must erase that bit, paying back the very
$k_{\mathrm{B}} T \ln 2$ it gained. The engine extracts $k_{\mathrm{B}} T \ln 2$ from a bit, and Landauer's
bound charges $k_{\mathrm{B}} T \ln 2$ to clear it: the books close, the cycle yields nothing, and the exact
equality of the two is the sharpest statement in physics that information is physical. A bit is worth
$k_{\mathrm{B}} T \ln 2$ --- to spend, and to clear.

The honest floor is a finite count obligation: the erased bit is a count, and $k_{\mathrm{B}} T \ln 2$ is its
price. The reading is that the irreversible step is not measurement but erasure. A record can be made for
free, but it cannot be unmade for free. The second law is paid at the moment a distinction is destroyed, not
when it is first drawn --- which is why the demon could sort all day for nothing and still never beat the law:
the bill was waiting at the bottom of the notebook.

\section{The Entropic Cost of Acceleration: temperature from a horizon}
\label{se:unruh}

Temperature, the chapter's third effect shows, can be made from geometry alone. Accelerate, and a horizon
forms --- a surface beyond which the records of the world can no longer reach you, sealed off by your own
motion. The records behind it are not destroyed but hidden, and their count is a temperature: an accelerating
observer reads a thermal bath whose warmth rises with the acceleration,
\begin{equation}
  T \propto a,
\end{equation}
the heat set by the entropy of the records the horizon has sealed away. Temperature here is not primary, and
not a property of any substance; it is the count of sealed pages the horizon hides.

This is the instrument meeting a horizon again, in a new guise. In Part I the butterfly gave it a horizon of
\emph{resolution} --- a depth past which its own floor swamped the signal. Here acceleration gives it a
horizon of \emph{access} --- a surface past which records are sealed from it. And the same idea returns at the
largest scale in Part V, where a black hole's event horizon seals records away and carries, by exactly this
reasoning, a temperature and an entropy of its own. The Unruh effect is the bench-scale rehearsal of the
black hole: a horizon, a sealed count, a temperature read off what can no longer be reached.

The honest floor is a smooth-shadow analogy: the discrete count of hidden records is the shadow of the
continuum temperature, and the bridge from the sealed-page count to $T \propto a$ is an analogy, named as one.
The reading is that temperature can be made from geometry alone. Accelerate, and a horizon seals part of your
own ledger away; the heat you then feel is the entropy of what you can no longer read --- temperature as the
price of a horizon, not a property of a substance.

\section{The Ideal Ledger Effect: pressure as reconciliation flux}
\label{se:ideal-ledger}

Pressure, too, turns out to be bookkeeping rather than substance. To the instrument, pressure is the rate at
which a boundary must reconcile the updates arriving from the interior: every molecule that reaches the wall
is a reconciliation the wall must perform, and the pressure is how fast those reconciliations come. When the
volume is large the requests are sparse and the wall is rarely troubled; when it is small they crowd the same
stretch of boundary and the wall is hammered. Pressure is the flux density of that reconciliation,
\begin{equation}
  P \propto \frac{n\,T}{V}, \qquad\text{that is,}\qquad PV = nT,
\end{equation}
the ideal gas law read as additive ledger counts --- $n$ threads each reconciling at a shared wall, the
temperature $T$ setting how vigorously each one arrives, the volume $V$ setting how much boundary they have to
share.

Pressure is one reading among many, and all of them flow from a single object the instrument writes down once
and differentiates forever: the partition function. List the states a system can occupy, each of energy $E_i$,
and sum a Boltzmann weight over them,
\begin{equation}
  Z = \sum_i e^{-\beta E_i}, \qquad \beta = \frac{1}{k_{\mathrm{B}} T},
\end{equation}
the weight $e^{-\beta E_i}$ measuring how accessible a state of energy $E_i$ is at temperature $T$ --- near one
for states well below $k_{\mathrm{B}} T$, vanishing for states far above. The probability of finding the system
in state $i$ is that weight normalized,
\begin{equation}
  p_i = \frac{e^{-\beta E_i}}{Z},
\end{equation}
the Boltzmann distribution. And from $Z$ the whole thermodynamics falls out by differentiation. Its logarithm
is the free energy,
\begin{equation}
  F = -k_{\mathrm{B}} T \ln Z,
\end{equation}
and every thermodynamic count the chapter has named is a derivative of that one sum over states: the average
energy is $\langle E\rangle = -\partial_\beta \ln Z$, the entropy is $S = -\partial_T F$, the pressure is $P =
-\partial_V F$. The instrument does not compute heat, entropy, and pressure separately; it writes down one sum
over the accessible states and reads each observable off as a slope. $Z$ is the ledger of the many, and
thermodynamics is its bookkeeping.

The honest floor is a finite count obligation: pressure is a count of reconciliation events per unit of
boundary, and $PV = nT$ balances those counts. The reading is that pressure is not a substance pushing on a
wall but the density of reconciliation the wall must carry. The ideal gas law is bookkeeping --- more threads,
a hotter gas, or a smaller wall each mean more reconciliations crowded together, and that crowding is the
pressure.

\section{The Thermostat Effect: equilibrium actively held}
\label{se:thermostat}

Equilibrium is the last surprise, and the gentlest: it is not a passive resting point but an actively
maintained one. An admissible ledger self-regulates around its low-strain states, and its updates do not
merely seek a stationary point --- they suppress deviations away from it. The stability is second-order. Where
the curvature of the strain is positive,
\begin{equation}
  \delta^2 \mathcal{I} > 0,
\end{equation}
a deviation decays and the ledger returns, a stable minimum; where $\delta^2 \mathcal{I} < 0$, the deviation
amplifies and the control loop fails. Equilibrium is the fixed point the feedback holds, not merely the place
the system happens to land.

The fixed point has a name from the partition function of the last section. Equilibrium is the state that
minimizes the free energy $F = -k_{\mathrm{B}} T \ln Z$ --- equivalently, the state of maximum entropy at fixed
energy, the most probable arrangement of the many. The first variation of $F$ vanishes there, $\delta F = 0$,
marking the extremum; the second variation is positive, $\delta^2 F > 0$, marking it a minimum rather than a
maximum --- which is exactly the $\delta^2 \mathcal{I} > 0$ above, the curvature that makes a deviation decay.
Equilibrium is not a place the system drifts to but the bottom of the free-energy basin, and the second law's
drift toward it is the system rolling down that basin while the feedback holds it at the floor.

The honest floor is a finite count obligation: stability is the sign of the second variation, a definite
verdict on whether a deviation grows or shrinks. The reading is that equilibrium is maintained, not merely
occupied. A thermostat does not so much find a stable state as hold one, suppressing drift by negative
feedback --- and the second law's drift toward equilibrium is itself an active reconciliation, the ledger
forever restoring the consistency a fluctuation disturbs, not a passive fall toward disorder.

\bigskip

The thread through all five is one claim: information is physical, and the instrument's own measuring is part
of the physics it measures. Sorting, erasing, accelerating, confining, and stabilizing each cost in the
currency of the second law, because each one makes, destroys, hides, crowds, or holds a count of distinctions
--- and that count is the entropy. The demon makes the price of a measurement visible; Landauer collects it at
erasure; a horizon turns a sealed count into heat; a wall turns crowded reconciliations into pressure; and a
thermostat spends a steady trickle to hold a fixed point against drift. Part III closes next with matter in
bulk, where these same costs harden into the properties of solids.
